\documentclass[10pt,a4paper]{amsart}
\usepackage[margin=19mm]{geometry}
\usepackage{amsmath,amssymb,mathtools}
\usepackage[scr=rsfs,cal=euler]{mathalfa}
\usepackage{xfrac}
\usepackage[unicode,hidelinks,bookmarks=false]{hyperref}
\newcommand{\ie}{i.e.\ }
\newcommand{\cf}{cf.\ }
\newcommand{\resp}{respectively\ }
\newcommand{\Subsection}{\subsection}
\providecommand{\nobreakdash}{\nobreak\hbox{-}}
\setlength{\emergencystretch}{2em}
\multlinegap0pt
\usepackage{bm}
\usepackage{bbm}
\usepackage{dsfont}
\usepackage{xspace}
\usepackage{cancel}
\usepackage{mathtools}
\usepackage{amsthm}
\makeatletter
\@ifundefined{thm}{\newtheorem{thm}{Theorem}}{}
\@ifundefined{lem}{\newtheorem{lem}[thm]{Lemma}}{}
\@ifundefined{prop}{\newtheorem{prop}[thm]{Proposition}}{}
\makeatother
\newcommand{\CAT}{\operatorname{CAT}}
\newcommand{\GH}{\operatorname{\mathrm{GH}}}
\newcommand{\N}{\operatorname{\mathbb{N}}}
\newcommand{\R}{\operatorname{\mathbb{R}}}
\newcommand{\Sph}{\operatorname{\mathbb{S}}}
\title[Asymptotic geometric regularity of CAT(0) spaces]{Asymptotic geometric regularity of CAT(0) spaces}
\author{Koichi Nagano}
\date{Source: arXiv:2602.20533v1 (CC BY 4.0)}
\begin{document}
\maketitle

\noindent\emph{Source and licence.} Asymptotic geometric regularity of CAT(0) spaces. Version arXiv:2602.20533v1 (CC BY 4.0), distributed under the Creative Commons Attribution 4.0 International licence, as recorded in the arXiv licence metadata for this exact version and shown on the abstract page https://arxiv.org/abs/2602.20533v1. Full rights notice: licenses/arxiv-2602.20533-CC-BY-4.0.txt.\par\smallskip

\noindent\textit{Editorial scope.} Counted unit: the Thm whose source label is \texttt{thm: grcat1g} in the article (source0.tex lines 1941--1950, source label \texttt{thm: grcat1g}), delivered together with the article's own complete original proof of it (source0.tex lines 1952--2114, one \texttt{proof} environment of 163 lines). No mathematics is written, repaired or re-derived: statement and proof are the article's own text line for line. Required context retained uncounted: lem: aperp (lines 960--981); lem: sphstrinf (lines 1135--1150); prop: fullbstr (lines 1834--1844). Every definition, display and label of those lines is retained unchanged. Omitted from this package and not redistributed: 1--959, 982--1134, 1151--1833, 1845--1940, 1951--1951, 2115--2288. Nothing inside a retained excerpt is omitted or altered: every retained body is delivered exactly as the article's lines are written, so no omission receipt of any kind accompanies this package. Citations: the retained excerpts carry no citation command at all, so no bibliography entry of the article is transcribed and no reference is redistributed. Reference adaptation: none was needed and none was made; every reference inside the delivered unit resolves inside it, so no unresolved reference or citation is produced. Printed slips kept literal: no printed word, symbol, hypothesis or number of the counted statement or of its proof was altered. Layout and class: the article's own class is \texttt{amsart} with packages including amsmath, amssymb, amsthm, eucal; the wrapper is the lane's pinned \texttt{amsart} editorial wrapper at 10pt on A4, which restates the theorem-like environments the retained text needs and adds the article's own macro definitions that the retained text needs (\texttt{\textbackslash CAT}, \texttt{\textbackslash GH}, \texttt{\textbackslash N}, \texttt{\textbackslash R}, \texttt{\textbackslash Sph}), with extra wrapper packages (bm, bbm, dsfont, xspace, cancel, mathtools). The delivered numbering is the unit's own. Assets: no retained span contains a figure environment, a graphics-inclusion command, a PDF-page inclusion, a table or a TikZ picture; no image file is copied into this package, the package declares no asset dependency and no image file is redistributed with it. Licence: version arXiv:2602.20533v1 of this article is distributed under the Creative Commons Attribution 4.0 International licence, as recorded in the arXiv licence metadata for this exact version and shown on the abstract page https://arxiv.org/abs/2602.20533v1. A full rights notice is delivered at licenses/arxiv-2602.20533-CC-BY-4.0.txt. Characters: every retained excerpt is pure ASCII, so the delivered engine, XeLaTeX with its default Latin Modern text font, reports no missing character.
\begin{lem}\label{lem: aperp}
For every $\epsilon \in (0,1)$,
and for every $n \in \N$,
there exists a sufficiently small $\delta \in (0,1)$ 
satisfying the following property:
Let $Z$ be a penetrable $\CAT(1)$ space with metric $d_Z$
of $\dim_{\mathrm{G}}Z \le n-1$
admitting an $(n,\delta)$-suspender
$(p_1,\dots,p_n)$.
If $z_1, z_2 \in Z$ satisfy 
\[
\left\vert d_Z(z_1,z_2) - \frac{\pi}{2} \right\vert < \delta,
\]
then we have
\[
\left\vert
\sum_{i=1}^n \cos d_Z \left( p_i, z_1 \right)
\cos d_Z \left( p_i, z_2 \right)
\right\vert
< \epsilon.
\]
\end{lem}

\begin{lem}\label{lem: sphstrinf}
Let $X$ be a $\CAT(0)$ space,
and let $(\xi_1, \dots, \xi_m)$ be an $(m,\delta)$-strainer at infinity
in $\partial_{\mathrm{T}}X$.
Then for every $x \in X$ the $m$-tuple
$((\xi_1)_x', \dots, (\xi_m)_x')$ of the directions at $x$
is an $(m,\delta)$-suspender in $\Sigma_xX$.
Moreover,
if $(\xi_1, \dots, \xi_m)$ and $(\eta_1, \dots, \eta_m)$ 
are mutually opposite $(m,\delta)$-strainers at infinity
in $\partial_{\mathrm{T}}X$,
then 
$\left( (\xi_1)_x', \dots, (\xi_m)_x' \right)$ and 
$\left( (\eta_1)_x', \dots, (\eta_m)_x' \right)$
are mutually opposite $(m,\delta)$-suspenders in $\Sigma_xX$.
\end{lem}

\begin{prop}\label{prop: fullbstr}
For every $\epsilon \in (0,1)$,
and for every $n \in \N$,
there exists 
$\delta \in (0,1)$ satisfying the following:
Let $X$ be a geodesically complete,
complete $\CAT(0)$ space of $\dim_{\mathrm{G}}X \le n$.
If $\varphi \colon X \to \R^n$
is a Busemann $(n,\delta)$-strainer map on $X$,
then it is a $(1+\epsilon)$-bi-Lipschitz homeomorphism.
\end{prop}

\begin{thm}
\label{thm: grcat1g}
For every $\epsilon \in (0,1)$,
and for every $n \in \N$,
there exists $\delta \in (0,1)$ such that
if a geodesically complete $\CAT(1)$ space $Z$
of $\dim_{\mathrm{G}} Z \le n-1$ satisfies
$d_{\GH} \left( Z, \Sph^{n-1} \right) < \delta$,
then $Z$ is $(1+\epsilon)$-bi-Lipschitz homeomorphic to $\Sph^{n-1}$.
\end{thm}

\begin{proof}
For $n \in \N$,
let $\delta \in (0,1)$ be small enough.
Let $Z$ be a geodesically complete $\CAT(1)$ space with metric $d_Z$
of $\dim_{\mathrm{G}}Z \le n-1$
satisfying $d_{\GH} \left( Z, \Sph^{n-1} \right) < \delta$.
Then the Euclidean cone $C_0(Z)$ over $Z$
is a geodesically complete $\CAT(0)$ space
satisfying $\dim_{\mathrm{G}}C_0(Z) \le n$.
Since $Z$ is isometric to $\partial_{\mathrm{T}}C_0(Z)$,
we have 
$d_{\GH} \left( \partial_{\mathrm{T}}C_0(Z), \Sph^{n-1} \right) < \delta$.
In this case,
there exists an $(n,10\delta)$-strainer at infinity
$\left( \xi_1, \dots, \xi_n \right)$.
For each $i \in \{ 1, \dots, n \}$,
take a ray $\gamma_i \colon [0,\infty) \to C_0(Z)$
with $\xi_i = \gamma_i(\infty)$.
Then the map $\varphi_0 \colon C_0(Z) \to \R^n$
defined by 
$\varphi_0 = \left( b_{\gamma_1}, \dots, b_{\gamma_n} \right)$
is a Busemann $(n,10\delta)$-strainer map.
Due to Proposition \ref{prop: fullbstr},
the map
$\varphi_0$ is a 
$\left( 1+\vartheta_n(\delta) \right)$-bi-Lipschitz homeomorphism.

Identify the unit metric sphere centered at the vertex in $C_0(Z)$ with 
$Z$,
and do the unit metric sphere centered at the origin in $\R^n$ with 
$\Sph^{n-1}$.
Define a bijective map $\varphi \colon Z \to \Sph^{n-1}$
by
\[
\varphi(z) := \frac{1}{\Vert \varphi_0(z) \Vert}.
\]
Take distinct points $z_1, z_2$ in $Z$.
If $z_1$ and $z_2$ satisfy $d_Z(z_1,z_2) \ge \pi/100$,
then we see 
\[
\left\vert
d_{\Sph^{n-1}} \left( \varphi(z_1), \varphi(z_2) \right) - 
d_Z \left( z_1, z_2 \right) 
\right\vert
< \left( 1+ \vartheta_n(\delta) \right) d_Z \left( z_1, z_2 \right),
\]
where $d_{\Sph^{n-1}}$ is the metric on $\Sph^{n-1}$.

Assume that 
$z_1$ and $z_2$ satisfy $d_Z(z_1,z_2) < \pi/100$.
Set $l := d_Z(z_1,z_2)$.
Let $\gamma \colon [0,l] \to Z$ be the geodesic
in $Z$ from $z_1$ to $z_2$.
Take $t \in [0,l)$.
By Lemma \ref{lem: sphstrinf},
the $n$-tuple 
$\left( (\xi_1)_{\gamma(t)}', \dots, (\xi_n)_{\gamma(t)}' \right)$
is an $(n,10\delta)$-strainer in $\Sigma_{\gamma(t)}C_0(Z)$.
Let $\xi_{\gamma(t)}$ be the direction $\left (D_t\gamma \right)(1)$
in $\Sigma_{\gamma(t)}C_0(Z)$.
Since $\varphi_0$ has Busemann function coordinates,
we have
\begin{align*}
\left( D_{\gamma(t)} \varphi_0 \right) \left( \xi_{\gamma(t)} \right)
&= - \left( \cos \angle_{\gamma(t)}
\left( \xi_1, \xi_{\gamma(t)} \right), \dots,
\cos \angle_{\gamma(t)} \left( \xi_n, \xi_{\gamma(t)} \right) \right), \\
\left( D_{\gamma(t)} \varphi_0 \right) \left( 0_{\gamma(t)}' \right)
&= - \left( \cos \angle_{\gamma(t)} \left( \xi_1, 0_{\gamma(t)}' \right), 
\dots,
\cos \angle_{\gamma(t)} \left( \xi_n, 0_{\gamma(t)}' \right) \right).
\end{align*}
Since $\varphi_0$ is a 
$\left( 1+\vartheta_n(\delta) \right)$-bi-Lipschitz homeomorphism,
we have
\[
\left\vert 
\left\Vert
\left( D_{\gamma(t)} \varphi_0 \right) \left( \xi_{\gamma(t)} \right) 
\right\Vert - 1
\right\vert < \vartheta_n(\delta),
\quad
\left\vert 
\left\Vert
\left( D_{\gamma(t)} \varphi_0 \right) \left( 0_{\gamma(t)}' \right) 
\right\Vert - 1
\right\vert < \vartheta_n(\delta).
\]
Applying Lemma \ref{lem: aperp} to $\Sigma_{\gamma(t)}C_0(Z)$,
we see
\[
\left\vert 
\angle_{\gamma(t)} \left(
\left( D_{\gamma(t)} \varphi_0 \right) \left( \xi_{\gamma(t)} \right),
\left( D_{\gamma(t)} \varphi_0 \right) \left( 0_{\gamma(t)}' \right) \right)
- \frac{\pi}{2} 
\right\vert < \vartheta_n(\delta);
\]
in other words,
the tangent vector 
$\left( D_{\gamma(t)} \varphi_0 \right) \left( \xi_{\gamma(t)} \right)$
of $\varphi_0 \circ \gamma$ at $t$
and the vector from $(\varphi_0 \circ \gamma)(t)$ to the origin $0$ 
are almost perpendicular.
Hence we see
\[
\left\vert
\frac{
\left\Vert \left( D_{\gamma(t)} \varphi \right) 
\left( \xi_{\gamma(t)} \right) \right\Vert}
{\left\Vert \left( D_{\gamma(t)} \varphi_0 \right) 
\left( \xi_{\gamma(t)} \right) \right\Vert}
-1
\right\vert < \vartheta_n(\delta).
\]
Therefore we obtain
\begin{multline*}
d_{\Sph^{n-1}} \left( \varphi(z_1), \varphi(z_2) \right)
\le \int_0^l \left\Vert \left( D_{\gamma(t)} \varphi \right) 
\left( \xi_{\gamma(t)} \right) \right\Vert dt \\
< \left( 1+\vartheta_n(\delta) \right) 
\int_0^l \left\Vert \left( D_{\gamma(t)} \varphi_0 \right) 
\left( \xi_{\gamma(t)} \right) \right\Vert dt
< \left( 1+\vartheta_n(\delta) \right) d_Z(z_1,z_2).
\end{multline*}

Set $\bar{l} := d_{\Sph^{n-1}} \left( \varphi(z_1), \varphi(z_2) \right)$.
Let $\bar{\gamma} \colon [0,\bar{l}] \to \Sph^{n-1}$ be 
the geodesic in $\Sph^{n-1}$
from $\varphi(z_1)$ to $\varphi(z_2)$, and
$\tilde{\gamma} \colon [0,l] \to \Sph^{n-1}$
the constant-speed reparametrization 
of $\bar{\gamma}$ defined by 
$\tilde{\gamma}(t) := \bar{\gamma} \left( (\bar{l}/l)t \right)$.
Since the tangent vector 
$\left( D_{\gamma(t)} \varphi_0 \right) \left( \xi_{\gamma(t)} \right)$
and the vector from $(\varphi_0 \circ \gamma)(t)$ to the origin $0$
are almost perpendicular,
we see
\[
\left\vert
\frac
{\left\Vert \left( D_{\gamma(t)} \varphi_0 \right) 
\left( \xi_{\gamma(t)} \right) \right\Vert}
{\left\Vert \left( D_t\tilde{\gamma} \right) (1) \right\Vert}
-1
\right\vert < \vartheta_n(\delta).
\]
Therefore we obtain
\begin{multline*}
d_{\Sph^{n-1}} \left( \varphi(z_1), \varphi(z_2) \right)
= \int_0^l \left\Vert \left( D_t\tilde{\gamma} \right) (1) \right\Vert dt \\
> \left( 1-\vartheta_n(\delta) \right) 
\int_0^l \left\Vert \left( D_{\gamma(t)} \varphi_0 \right) 
\left( \xi_{\gamma(t)} \right) \right\Vert dt
= \left( 1-\vartheta_n(\delta) \right) d_Z(z_1,z_2).
\end{multline*}
Thus $\varphi$ is a 
$\left( 1+\vartheta_n(\delta) \right)$-bi-Lipschitz homeomorphism.

In this way, 
we have completed the proof of Theorem \ref{thm: grcat1g}.
\end{proof}

\end{document}
